\begin{lemma}
\label{lemma-construct-bijective-orbits}
Let $\mathcal{F}$ be a predeformation category satisfying
(S2) which has a versal formal object. Then its minimal versal
formal object satisfies (\ref{equation-bijective-orbits}).
\end{lemma}

\begin{proof}
Let $\xi$ be a minimal versal formal object for $\mathcal{F}$, see
Lemma \ref{lemma-minimal-versal}.
Say $\xi$ lies over $R \in \Ob(\widehat{\mathcal{C}}_\Lambda)$.
In order to parse (\ref{equation-bijective-orbits}) we point out
that $T\mathcal{F}$ has a natural $k$-vector space structure
(see
Lemma \ref{lemma-tangent-space-vector-space}),
that $d\underline{\xi} : \text{Der}_\Lambda(R, k) \to T\mathcal{F}$
is linear (see
Lemma \ref{lemma-k-linear-differential}),
and that the action of $\text{Der}_\Lambda(k, k)$ is
given by addition (see
Lemma \ref{lemma-action-linear}).
Consider the diagram
$$
\xymatrix{
& \Hom_k(\mathfrak m_R/\mathfrak m_R^2, k) \\
K \ar[r] & \text{Der}_\Lambda(R, k) \ar[r]^{d\underline{\xi}} \ar[u] &
T\mathcal{F} \\
& \text{Der}_\Lambda(k, k) \ar[u] \ar[ru]
}
$$
The vector space $K$ is the kernel of $d\underline{\xi}$.
Note that the middle column is exact in the middle as it is dual to the
sequence (\ref{equation-sequence}). If (\ref{equation-bijective-orbits})
fails, then we can find a nonzero element $D \in K$ which
does not map to zero in $\Hom_k(\mathfrak m_R/\mathfrak m_R^2, k)$.
This means there exists an $t \in \mathfrak m_R$ such that
$D(t) = 1$. Set $R' = \{a \in R \mid D(a) = 0\}$. As $D$ is a derivation
this is a subring of $R$. Since $D(t) = 1$ we see that $R' \to k$
is surjective (compare with the proof of
Lemma \ref{lemma-essential-surjection}).
Note that $\mathfrak m_{R'} = \Ker(D : \mathfrak m_R \to k)$
is an ideal of $R$ and $\mathfrak m_R^2 \subset \mathfrak m_{R'}$. Hence
$$
\mathfrak m_R/\mathfrak m_R^2 =
\mathfrak m_{R'}/\mathfrak m_R^2 + k\overline{t}
$$
which implies that the map
$$
R'/\mathfrak m_R^2 \times_k k[\epsilon] \to R/\mathfrak m_R^2
$$
sending $\epsilon$ to $\overline{t}$ is an isomorphism. In particular
there is a map $R/\mathfrak m_R^2 \to R'/\mathfrak m_R^2$.

\medskip\noindent
Let $\xi \to y$ be a morphism lying over $R \to R/\mathfrak m_R^2$.
Let $y \to x$ be a morphism lying over
$R/\mathfrak m_R^2 \to R'/\mathfrak m_R^2$.
Let $y \to x_\epsilon$ be a morphism lying over
$R/\mathfrak m_R^2 \to k[\epsilon]$. Let $x_0$ be the unique (up to unique
isomorphism) object of $\mathcal{F}$ over $k$.
By the axioms of a category cofibred in groupoids we obtain
a commutative diagram
$$
\vcenter{
\xymatrix{
y \ar[r] \ar[d] & x_\epsilon \ar[d] \\
x \ar[r]   & x_0
}
}
\quad\text{lying over}\quad
\vcenter{
\xymatrix{
R'/\mathfrak m_R^2 \times_k k[\epsilon] \ar[r] \ar[d] & k[\epsilon] \ar[d] \\
R'/\mathfrak m_R^2 \ar[r] & k.
}
}
$$
Because $D \in K$ we see that $x_\epsilon$ is isomorphic
to $0 \in \mathcal{F}(k[\epsilon])$, i.e., $x_\epsilon$ is
the pushforward of $x_0$ via $k \to k[\epsilon], a \mapsto a$.
Hence by
Lemma \ref{lemma-lifting-section}
we see that there exists a morphism $x \to y$. Since
$\text{length}_\Lambda(R'/\mathfrak m_R^2) <
\text{length}_\Lambda(R/\mathfrak m_R^2)$
the corresponding ring map $R'/\mathfrak m_R^2 \to R/\mathfrak m_R^2$
is not surjective. This contradicts the minimality of
$\xi/R$, see part (1) of
Lemma \ref{lemma-smallest-where-descends-versal}.
This contradiction shows that such a $D$ cannot exist, hence
we win.
\end{proof}